\section{Statistical targets and guarantees}
\label{supp:statistics}

\subsection{Fixed regions and independent evaluation}
A region is a fixed measurable event for an image. A validated semantic
classifier can define it, or it can be fitted in feature space using the
reference-fit pool $R_{\rm fit}$.
Conditional on this fit, its probability is a fixed population quantity.
For independently sampled evaluation rows, each region count has a binomial
marginal distribution. The regions, representations, and configurations can be
dependent; the union-bound construction does not require independence between
their counts. Independence of the evaluation rows from the fitted regions is
essential. Fitting a clustering model and evaluating its reference occupancy on
the same rows does not meet this condition.

When semantic labels are used, the same labeling function, confidence threshold,
and rejection category are applied to real and generated images. The estimated
probabilities then concern those declared classifier regions. Accuracy measured
on held-out real images validates that domain, but does not establish equal
accuracy on every generated-image distribution. We therefore retain rejected
images in the denominator, report feature fidelity, and include uncurated
qualitative checks. A full feature-space partition has no unassigned mass;
this does not make a fitted Gaussian mixture's latent-component error zero.

\subsection{Exact regional mass intervals}
For $k$ successes in $n$ independent Bernoulli trials and one-tail error
$\varepsilon$, let $L$ and $U$ be the lower and upper
Clopper--Pearson endpoints~\cite{clopper1934confidence}.
Here $\operatorname{Beta}^{-1}(u;a,b)$ is the $u$-quantile of a beta
law with shape parameters $a,b$.
\begin{align}
L(k,n;\varepsilon)&=\begin{cases}0,&k=0,\\
\operatorname{Beta}^{-1}(\varepsilon;k,n-k+1),&k>0,
\end{cases}\label{eq:cp_lower}\\
U(k,n;\varepsilon)&=\begin{cases}1,&k=n,\\
\operatorname{Beta}^{-1}(1-\varepsilon;k+1,n-k),&k<n.
\end{cases}\label{eq:cp_upper}
\end{align}
Each endpoint has noncoverage probability at most $\varepsilon$.
Let $p,q$ be the reference and candidate masses of a fixed region,
with bounds denoted by superscripts $-$ and $+$.
For $M$ prespecified region/configuration/condition/representation comparisons
and family error budget $\alpha_{\rm occ}$, set $\varepsilon=\alpha_{\rm occ}/(4M)$.
A union bound over both endpoints of both
distributions gives an event of probability at least $1-\alpha_{\rm occ}$ on
which every $p$ and $q$ is in its interval. On that event, for $p>0$,
\begin{equation}
q^-/p^+\leq q/p\leq q^+/p^-.
\label{eq:supp_occupancy}
\end{equation}
The right endpoint is defined to be infinite when $p^-=0$.
Thresholding these simultaneous intervals remains valid when several deficits
are reported together or when a fixed-sample comparison is stopped after its
first passing candidate. The guarantee does not license optional stopping
within an adaptively sized arm without additional sequential inference.

An absolute mass target uses the same endpoints to bound
$q-p\in[q^- -p^+,q^+-p^-]$. For a prespecified tolerance $\delta$,
intervals wholly above $\delta$ certify excess, wholly below $-\delta$
certify depletion, and wholly inside $[-\delta,\delta]$ certify equivalence.
Other cases are inconclusive. Section~\ref{sec:support-mass-balance} gives
the shared-reference allocation and persistence rule used in the signed audit.

\paragraph{A useful zero-count planning relation.}
With zero observed counts among $n_Q$ candidate draws, the upper
probability bound is $q^+=1-\varepsilon^{1/n_Q}$.
For an eligible region and a ratio threshold $\rho$ with
$0<\rho p^-<1$, certifying $q/p<\rho$ requires
\begin{equation}
n_Q>\frac{\log\varepsilon}{\log(1-\rho p^-)}.
\label{eq:zerocount}
\end{equation}
For small $\rho p^-$, this is approximately
$\log(4M/\alpha_{\rm occ})/(\rho p^-)$.
Reference uncertainty also matters: a small or ineligible $p^-$ can make the
test inconclusive despite zero generated counts. This relation explains why
the method reports a resolution and sample requirement rather than treating
zero observed categories as an unconditional claim of deletion.

\subsection{Support calibration and excess-tail uncertainty}
Let $S$ be the support score, fitted independently of the calibration
pool $R_{\rm cal}$. For $n_{\rm cal}$ reference calibration samples,
set $\tau$ to the score at rank $\lceil(n_{\rm cal}+1)(1-\eta)\rceil$,
in ascending order. The target tail level is $\eta=0.05$ by default.
Exchangeability with a new reference draw gives marginal split-conformal
tail control~\cite{vovk2013conditional}. If the rank exceeds the
calibration size, set $\tau=\infty$ and mark calibration uninformative.
For reference and candidate laws $P,Q$, independent evaluation rows
estimate tail probabilities $p_{\rm tail}=P(S>\tau)$ and
$q_{\rm tail}=Q(S>\tau)$.
Four exact binomial endpoints yield
\begin{equation}
q_{\rm tail}^- - p_{\rm tail}^+
\leq q_{\rm tail}-p_{\rm tail}
\leq q_{\rm tail}^+-p_{\rm tail}^-.
\label{eq:supp_tail_excess}
\end{equation}
These bounds account for reference-tail uncertainty rather than assuming its
observed rate is exactly five percent. Multiplicity is allocated across all
prespecified configurations, conditions, and representations. The feature-only
default allocates 0.025 each to occupancy and support. The fresh digit
routing and signed-balance studies instead allocate 0.025 to mass, 0.0125
to paired fidelity, and 0.0125 to support; each declared exact family totals
0.05. The original digit experiment separately allocates 0.025 to semantic
occupancy and 0.0125 each to feature occupancy and support; its fidelity
intervals use a three-training-seed Student-$t$ approximation.
This joint exact statement does not include the approximate spread intervals
or arbitrary external perceptual metrics.

\subsection{Conditional spread intervals}
For a feature sample $x_1,\ldots,x_n$ with mean $\bar x$ and sample
covariance trace $\hat v$, let $d_i=\|x_i-\bar x\|^2$.
Subscripts $P,Q$ index the reference and candidate samples, including
their sizes $n_P,n_Q$ and squared-distance vectors $d_P,d_Q$.
The estimated ratio is $\hat r=\hat v_Q/\hat v_P$.
The two-sample influence-function approximation
below assumes independent reference and candidate rows. For shared-noise
pairs in the crossed MNIST probes (Sec.~\ref{supp:mnist_routing}), we instead
estimate the variance of the paired influence-function differences, accounting
for covariance even when the two members enter different regions. For independent
samples, the log variance ratio uses
\begin{equation}
\widehat{\operatorname{se}}^2(\log\hat r)=
\frac{\widehat{\operatorname{Var}}(d_P)}{n_P\hat v_P^2}
+\frac{\widehat{\operatorname{Var}}(d_Q)}{n_Q\hat v_Q^2}.
\label{eq:supp_spread_se}
\end{equation}
Here $\widehat{\operatorname{se}}$ denotes estimated standard error.
The log-normal interval uses the prespecified Bonferroni level for the spread
family. It is approximate, especially in small, heavy-tailed cells. The code
also provides an independent-row percentile bootstrap. Cells with fewer than
30 samples in either distribution are unresolved. Null simulations and the
reported sensitivity results assess the actual approximation in the tested
settings; they do not establish distribution-free coverage.

\subsection{Prompt-level inference}
For prompt $c$, let $v_{Pc}$ and $v_{Qc}$ denote within-prompt trace covariance
estimates. A bar denotes the average over prompts. The primary ratio is
$\bar v_Q/\bar v_P$, rather than the average of $v_{Qc}/v_{Pc}$.
Matched prompts are resampled as units in the bootstrap, preserving all their
seeds and arm pairings. The same principle applies to alignment, preference,
and per-prompt nearest-neighbor metrics. Calibration of a population tail uses
one score-blind seed per independent calibration prompt. Pooling all images
from repeated prompts into an independent-row binomial calculation would
overstate the evidence. With only a few seeds per prompt, semantic occupancy
tests generally lack power; they do not establish absence of prompt-specific
depletion.

\section{Controlled comparisons identify useful interventions}
\label{supp:causality}
Let $Z\sim\mathcal N(0,I_d)$ be a standard Gaussian in $\mathbb R^d$,
where $I_d$ is the identity, and let $\gamma\in(0,1)$.
In one world, the learned resolved generator outputs $Z$ and an approximate
sampler outputs $\gamma Z$. In another, the learned generator outputs
$\gamma Z$ and its sampler is exact. Both observed distributions are
$\mathcal N(0,\gamma^2 I_d)$, with full support $\mathbb R^d$.
Every final-sample statistic has the same law in the two worlds. Analogously,
a sampler that reweights a two-component teacher can have the same output
mixture as a model trained with those weights and sampled exactly. Therefore,
regional depletion or contraction alone cannot determine whether a sampling,
conditioning, or training intervention will help. The two-reference procedure
adds information and measures the response to controlled changes. These elementary identification limits motivate the intervention comparisons.

For a feature vector $X$ with mean $\mu$ and regional index $J$, let
$\mu_J$ be its mean conditional on $J$. The total-covariance identity
follows by writing
$X-\mu=(X-\mu_J)+(\mu_J-\mu)$ and conditioning on $J$. The cross terms vanish
because $\mathbb E[X-\mu_J\mid J]=0$. Taking traces separates the contributions
of within-region covariance and between-region means, each weighted by regional
mass. A rotation of an anisotropic conditional distribution preserves its trace
but may change a reference-support tail. This motivates the rotation control
and the distinction between support departure and expansion.

